\documentclass[11pt,a4paper]{article}
\usepackage[utf8]{inputenc}
\usepackage[italian]{babel}
\usepackage{geometry}
\geometry{a4paper, margin=2.5cm}
\usepackage{enumitem}
\usepackage{amsmath}

\title{\textbf{Proposta sviluppo modulo AI} \\ \Large Radio generata da prompt}
\author{Gabriele Rubboli Petroselli}
\date{10/03/2026}

\begin{document}

\maketitle
\vspace{1.5cm}

\section{Obiettivo del modulo AI}

Il modulo AI ha lo scopo di generare automaticamente una trasmissione radio a partire da un prompt testuale dell'utente.

\paragraph{Input esempio:}
``Musica chill per studiare la sera''

\paragraph{Output del sistema:}
\begin{enumerate}
    \item Scaletta di canzoni coerente con il prompt
    \item Testo dello speaker radio tra una canzone e l'altra
\end{enumerate}

\paragraph{Esempio output:}
\textit{Playlist:}
\begin{enumerate}
    \item Song A -- Artist X
    \item Song B -- Artist Y
    \item Song C -- Artist Z
\end{enumerate}

\textit{Speaker text:}

``Benvenuti alla nostra selezione musicale serale. Iniziamo con\dots''

\vspace{0.5cm}
Il sistema restituirà l'output in forma strutturata (JSON/API) per essere integrato nel sistema principale che gestirà:
\begin{itemize}
    \item player audio
    \item generazione della voce dello speaker
    \item streaming radio
\end{itemize}

\vspace{0.5cm}
\hrulefill
\vspace{0.5cm}

\section*{Ambito del lavoro}

Il lavoro riguarda solo il modulo AI.

\paragraph{Incluso:}
\begin{itemize}
    \item interpretazione del prompt utente
    \item selezione delle canzoni
    \item generazione della scaletta
    \item generazione del testo dello speaker
    \item output strutturato per integrazione backend
\end{itemize}

\paragraph{Non incluso:}
\begin{itemize}
    \item sviluppo frontend
    \item player audio
    \item streaming radio
    \item gestione utenti
    \item infrastruttura completa dell'app
\end{itemize}

\vspace{0.5cm}
\hrulefill
\vspace{0.5cm}

\section{Opzione 1 -- Prototipo semplice}

\paragraph{Descrizione}
Versione minima del sistema.

Il prompt viene interpretato da un LLM che estrae informazioni come:
\begin{itemize}
    \item genere
    \item mood
    \item livello di energia
\end{itemize}

Queste informazioni vengono utilizzate per interrogare direttamente Spotify Recommendations API.

\paragraph{Pipeline:}
\begin{center}
Prompt utente \\
$\downarrow$ \\
LLM interpreta il prompt \\
$\downarrow$ \\
Conversione in target musicali \\
$\downarrow$ \\
Spotify Recommendations API \\
$\downarrow$ \\
Selezione playlist \\
$\downarrow$ \\
LLM genera testo speaker
\end{center}

\paragraph{Esempio chiamata Spotify:}
\begin{verbatim}
GET /recommendations
seed_genres=ambient,chill
target_energy=0.3
target_tempo=80
target_valence=0.5
\end{verbatim}

Spotify restituisce direttamente una lista di brani.

Il sistema utilizza quasi direttamente questi risultati.

\paragraph{Vantaggi}
\begin{itemize}
    \item sviluppo molto rapido
    \item costo basso
    \item sufficiente per demo o prototipo
\end{itemize}

\paragraph{Svantaggi}
\begin{itemize}
    \item controllo limitato sulla qualità della playlist
    \item forte dipendenza dagli algoritmi Spotify
    \item possibili ripetizioni di artista
\end{itemize}

\vspace{0.3cm}
\noindent \textbf{Tempo stimato:} 4 -- 7 giorni \\[0.2cm]
\noindent \textbf{Prezzo stimato:} €400 -- €700

\vspace{0.5cm}
\hrulefill
\vspace{0.5cm}

\section{Opzione 2 -- MVP con ranking musicale}

\paragraph{Descrizione}
Versione più robusta.

Spotify viene utilizzato per trovare candidate songs, ma la selezione finale viene fatta dal sistema AI.

\paragraph{Pipeline:}
\begin{center}
Prompt utente \\
$\downarrow$ \\
LLM interpreta il prompt \\
$\downarrow$ \\
Conversione in caratteristiche musicali target \\
$\downarrow$ \\
Spotify Recommendations / Search \\
$\downarrow$ \\
Pool di candidate songs \\
$\downarrow$ \\
Richiesta audio features delle candidate \\
$\downarrow$ \\
Ranking delle canzoni \\
$\downarrow$ \\
Playlist finale \\
$\downarrow$ \\
LLM genera speaker text
\end{center}

In questa versione il sistema utilizza le audio features reali delle canzoni (energy, tempo, valence ecc.) per calcolare quanto ogni brano è vicino al prompt.

Il ranking considera:
\begin{itemize}
    \item distanza dalle caratteristiche target
    \item popolarità del brano
    \item varietà artisti
    \item coerenza musicale
\end{itemize}

\paragraph{Vantaggi}
\begin{itemize}
    \item playlist più coerenti
    \item maggiore controllo sulla selezione musicale
    \item meno ripetizioni
\end{itemize}

\paragraph{Svantaggi}
\begin{itemize}
    \item sviluppo più lungo
    \item maggiore logica di ranking
\end{itemize}

\vspace{0.3cm}
\noindent \textbf{Tempo stimato:} 8 -- 12 giorni \\[0.2cm]
\noindent \textbf{Prezzo stimato:} €800 -- €1400

\vspace{0.5cm}
\hrulefill
\vspace{0.5cm}

\section{Opzione 3 -- Sistema avanzato con playlist planning}

\paragraph{Descrizione}
Questa versione migliora la qualità della trasmissione radio, non solo della playlist.

Oltre alla selezione delle canzoni, il sistema organizza la sequenza per creare una trasmissione più naturale.

\paragraph{Pipeline:}
\begin{center}
Prompt utente \\
$\downarrow$ \\
LLM interpreta il prompt \\
$\downarrow$ \\
Candidate songs via Spotify \\
$\downarrow$ \\
Ranking musicale \\
$\downarrow$ \\
Playlist planning \\
$\downarrow$ \\
Generazione speaker script
\end{center}

Il sistema ottimizza:
\begin{itemize}
    \item progressione dell'energia
    \item varietà degli artisti
    \item alternanza dei generi
    \item transizioni musicali
\end{itemize}

\paragraph{Vantaggi}
\begin{itemize}
    \item playlist più naturali
    \item esperienza radio più credibile
    \item migliore gestione delle transizioni musicali
\end{itemize}

\vspace{0.3cm}
\noindent \textbf{Tempo stimato:} 2 -- 3 settimane \\[0.2cm]
\noindent \textbf{Prezzo stimato:} €1400 -- €2200

\vspace{0.5cm}
\hrulefill
\vspace{0.5cm}

\section{Opzione 4 -- Sistema con Semantic Search (Embeddings)}

\paragraph{Descrizione}
Questa è l'architettura più avanzata.

Il sistema utilizza un catalogo musicale locale con embeddings semantici, che permette di trovare canzoni simili al significato del prompt.

In questo caso non si dipende solo da Spotify per trovare i brani.

\paragraph{Pipeline:}
\begin{center}
Prompt utente \\
$\downarrow$ \\
Embedding del prompt \\
$\downarrow$ \\
Ricerca semantica nel catalogo musicale \\
$\downarrow$ \\
Candidate songs semanticamente simili \\
$\downarrow$ \\
Ranking con audio features \\
$\downarrow$ \\
Playlist finale \\
$\downarrow$ \\
LLM genera speaker script
\end{center}

Il catalogo contiene per ogni canzone:
\begin{itemize}
    \item titolo
    \item artista
    \item generi
    \item anno
    \item audio features
    \item descrizione testuale
    \item embedding vettoriale
\end{itemize}

La ricerca avviene tramite vector database.

Questo permette di interpretare meglio prompt complessi, ad esempio:

``musica per guidare di notte sotto la pioggia''

\paragraph{Vantaggi}
\begin{itemize}
    \item maggiore comprensione semantica del prompt
    \item maggiore controllo sul catalogo musicale
    \item risultati più creativi e vari
\end{itemize}

\paragraph{Svantaggi}
\begin{itemize}
    \item sviluppo molto più complesso
    \item necessità di costruire un catalogo musicale locale
    \item maggiore lavoro di manutenzione dati
\end{itemize}

\vspace{0.3cm}
\noindent \textbf{Tempo stimato:} 3 -- 5 settimane \\[0.2cm]
\noindent \textbf{Prezzo stimato:} €2000 -- €3500

\end{document}